\documentclass{beamer}

% Theme Selection
\usetheme{Madrid}
\usecolortheme{beaver} % Red/Gray theme, professional for AP Econ

% Packages
\usepackage{graphicx}
\usepackage{amsmath}
\usepackage{tikz}
\usepackage{booktabs}

% Title Information
\title{Unit 8: Market Failure and the Role of Government}
\subtitle{Externalities, Public Goods, and Income Distribution}
\author{Doğukan Günkördü}
\date{}

\begin{document}

%---------------------------------------------------------
\begin{frame}
    \titlepage
\end{frame}

%---------------------------------------------------------
\begin{frame}{Unit Overview}
    \tableofcontents
\end{frame}

%---------------------------------------------------------
\section{Social Efficiency and Externalities}

\begin{frame}{Marginal Social Benefit vs. Marginal Social Cost}
    \begin{block}{Optimal Amount of Production}
        Social efficiency is achieved when all costs and benefits are taken into consideration.
        \[ \text{MSB} = \text{MSC} \]
    \end{block}

    \begin{itemize}
        \item \textbf{Example (Recycling):} Most people recycle because the marginal social benefit (MSB) outweighs the marginal social cost (MSC).
        \item \textbf{Example (Water Conservation):} Fewer people are willing to forgo a shower 3 times a week. At that point, the personal sacrifice (social cost) outweighs the benefit.
    \end{itemize}
    
    \vspace{0.5cm}
    \textbf{The Sweet Spot:}
    Society wants to produce where MSB = MSC. 
    \begin{itemize}
        \item If $MSB > MSC$: Society wants \textbf{more} built (Underproduction).
        \item If $MSC > MSB$: Society wants \textbf{fewer} built (Overproduction).
    \end{itemize}
\end{frame}

%---------------------------------------------------------
\begin{frame}{Defining Externalities}
    Externalities arise when producer or consumer actions result in spillover costs or benefits to \textbf{third parties}.

    \begin{columns}
        \begin{column}{0.5\textwidth}
            \textbf{Consumption Externalities}
            \begin{itemize}
                \item Arise from consumer activity.
                \item \textbf{Negative:} Cigarette smoke (secondhand smoke).
                \item \textbf{Positive:} Vaccines (herd immunity).
            \end{itemize}
        \end{column}
        \begin{column}{0.5\textwidth}
            \textbf{Production Externalities}
            \begin{itemize}
                \item Arise from producer activity.
                \item \textbf{Negative:} A firm dumping chemicals in a river.
                \item \textbf{Positive:} Research \& Development or subsidizing education.
            \end{itemize}
        \end{column}
    \end{columns}
\end{frame}

%---------------------------------------------------------
\begin{frame}{Positive Externalities}
    \begin{alertblock}{The Problem}
        The market fails to realize all benefits for society.
        \[ \text{MSB} > \text{MSC} \text{ at the market quantity.} \]
    \end{alertblock}

    \begin{itemize}
        \item \textbf{Example:} Higher education.
        \item \textbf{Spillover Benefits:} Lower crime rates, more productive economy, increased tax revenues.
        \item \textbf{Market Failure:} The market \textbf{underproduces} the good ($Q_{mkt} < Q_{social}$).
        \item \textbf{Deadweight Loss:} Exists at the market quantity because socially desirable trades are not happening.
    \end{itemize}
\end{frame}

%---------------------------------------------------------
\begin{frame}{Visualizing Positive Externalities}
    \textit{(Refer to Figure 11.2 and 11.3 in textbook)}

    \begin{itemize}
        \item \textbf{Production:} Two supply curves. MSC is below MPC (or MSB is higher than MPB).
        \item \textbf{Deadweight Loss Tip:} Imagine an arrow pointing to the optimal output where ``the socials meet'' ($MSB = MSC$). The base of the arrow is at the market equilibrium.
    \end{itemize}

    \begin{block}{The Government Solution: Subsidies}
        To correct underproduction, the government provides a \textbf{Per-Unit Subsidy}.
        \begin{itemize}
            \item Subsidy amount = Marginal External Benefit (MEB).
            \item Lump-sum subsidies do \textbf{not} change output efficiently; it must be per-unit.
        \end{itemize}
    \end{block}
\end{frame}

%---------------------------------------------------------
\begin{frame}{Negative Externalities}
    \begin{alertblock}{The Problem}
        The market produces at a cost greater to society than the firm pays.
        \[ \text{MSC} > \text{MSB} \text{ at the market quantity.} \]
    \end{alertblock}

    \begin{itemize}
        \item \textbf{Example:} Pollution, secondhand smoke, airplane noise.
        \item \textbf{Market Failure:} The market \textbf{overproduces} the good ($Q_{mkt} > Q_{social}$).
        \item \textbf{Analysis:} The firm ignores the external cost, leading to an inefficiently high quantity and low price.
    \end{itemize}
\end{frame}

%---------------------------------------------------------
\begin{frame}{Visualizing Negative Externalities}
    \textit{(Refer to Figure 11.4 and 11.5 in textbook)}
    
    \begin{itemize}
        \item \textbf{Graphing:} Supply (MPC) represents private costs. MSC is \textbf{higher} than MPC (shifted left/up) because it includes the external cost.
        \item \textbf{Deadweight Loss:} Points toward the social optimum (leftward).
    \end{itemize}

    \begin{block}{The Government Solution: Taxes}
        To correct overproduction, the government imposes a \textbf{Per-Unit Tax}.
        \begin{itemize}
            \item Tax amount = Marginal External Cost (MEC).
            \item This internalizes the externality, shifting the private cost curve to align with the social cost curve.
        \end{itemize}
    \end{block}
\end{frame}

%---------------------------------------------------------
\begin{frame}{Private Solutions: The Coase Theorem}
    \textbf{Question:} Is government intervention always necessary?
    
    \vspace{0.5cm}
    
    \textbf{The Coase Theorem (Ronald Coase):}
    \begin{itemize}
        \item Private parties can solve externality issues on their own via bargaining.
        \item \textbf{Condition:} Property rights must be clearly defined and transaction costs must be low.
        \item \textbf{Examples:}
        \begin{itemize}
            \item A polluting firm pays neighbors to compensate for harm.
            \item An airline buys soundproof windows for homes near an airport.
        \end{itemize}
    \end{itemize}
\end{frame}

%---------------------------------------------------------
\section{Public Goods}

\begin{frame}{Public vs. Private Goods}
    \begin{columns}
        \begin{column}{0.5\textwidth}
            \begin{block}{Private Goods}
                \begin{itemize}
                    \item \textbf{Excludable:} You can be prevented from using it if you don't pay.
                    \item \textbf{Rival:} One person's consumption reduces availability for others.
                    \item \textit{Example:} A ticket to a sold-out baseball game.
                \end{itemize}
            \end{block}
        \end{column}
        \begin{column}{0.5\textwidth}
            \begin{block}{Public Goods}
                \begin{itemize}
                    \item \textbf{Non-Excludable:} Cannot stop non-payers from benefiting.
                    \item \textbf{Non-Rival:} One person's use does not diminish another's.
                    \item \textit{Example:} National defense, clean air, lighthouses.
                \end{itemize}
            \end{block}
        \end{column}
    \end{columns}
\end{frame}

%---------------------------------------------------------
\begin{frame}{The Free-Rider Problem}
    \textbf{The Dilemma:}
    Since public goods are non-excludable, people can benefit without paying.
    
    \vspace{0.5cm}
    \textbf{Consequence:}
    \begin{itemize}
        \item If everyone free-rides, the market demand drops to zero.
        \item Private firms cannot make a profit, so they will \textbf{not supply} the good.
        \item \textbf{Solution:} The government must provide the good and fund it via taxes (e.g., taxpayer-funded fireworks shows).
    \end{itemize}
\end{frame}

%---------------------------------------------------------
\section{Income Distribution}

\begin{frame}{Measuring Inequality: The Lorenz Curve}
    Economics uses the Lorenz curve to visualize income distribution.
    
    \begin{itemize}
        \item \textbf{45-Degree Line:} Represents perfect income equality (50\% of families get 50\% of income).
        \item \textbf{Lorenz Curve:} The actual distribution (shaped like a banana).
    \end{itemize}

    \vspace{0.5cm}
    \begin{center}
        \textit{"The bigger the gap between the Lorenz curve and the perfect equality line (the bigger the banana), the more inequality exists."}
    \end{center}
\end{frame}

%---------------------------------------------------------
\begin{frame}{Measuring Inequality: The Gini Coefficient}
    \textbf{Gini Coefficient:} A mathematical ratio derived from the Lorenz Curve.
    
    \[ \text{Gini} = \frac{\text{Area between Line of Equality \& Lorenz Curve}}{\text{Total Area under Line of Equality}} \]

    \begin{itemize}
        \item \textbf{Range:} 0 to 1.
        \item \textbf{0:} Perfect Equality.
        \item \textbf{1:} Perfect Inequality (one person has all the income).
        \item \textit{Comparison:} The US has a higher Gini coefficient (more inequality) than France, but lower than Brazil.
    \end{itemize}
\end{frame}

%---------------------------------------------------------
\begin{frame}{Types of Taxes}
    Governments use taxes to redistribute income.
    
    \begin{enumerate}
        \item \textbf{Progressive Tax:}
        \begin{itemize}
            \item Tax rate \textbf{increases} as income increases.
            \item \textit{Example:} U.S. Federal Income Tax.
            \item Effect: Reduces income inequality.
        \end{itemize}
        
        \item \textbf{Proportional (Flat) Tax:}
        \begin{itemize}
            \item Tax rate stays the \textbf{same} regardless of income.
            \item Everyone pays the same percentage (e.g., 10\%).
        \end{itemize}
        
        \item \textbf{Regressive Tax:}
        \begin{itemize}
            \item Average tax burden \textbf{decreases} as percentage of income as income rises.
            \item \textit{Example:} Sales tax. (Low-income earners spend a higher \% of their total income on consumption than rich people).
        \end{itemize}
    \end{enumerate}
\end{frame}

%---------------------------------------------------------
\begin{frame}{Summary Formulas (Unit 8)}
    \begin{block}{Key Formulas}
        \begin{itemize}
            \item \textbf{Allocative Efficiency:} $MSB = MSC$
            \item \textbf{Social Benefits:} $MSB = MPB + MEB$
            \item \textbf{Social Costs:} $MSC = MPC + MEC$
        \end{itemize}
    \end{block}
    
    \vspace{0.5cm}
    \textbf{Identifying Externalities:}
    \begin{itemize}
        \item If $MSB > MSC$: Positive Externality (Underproduction).
        \item If $MSC > MSB$: Negative Externality (Overproduction).
    \end{itemize}
\end{frame}

%---------------------------------------------------------
\begin{frame}{End of Unit 8}
    \centering
    \Large
    \textbf{Review Questions?}
    \vspace{1cm}
    
    \small
    \textit{Remember: Positive externalities imply "Too little produced," Negative externalities imply "Too much produced."}
\end{frame}

\end{document}
